\let\negmedspace\undefined
\let\negthickspace\undefined
\documentclass[journal]{IEEEtran}
\usepackage[a5paper, margin=10mm, onecolumn]{geometry}
%\usepackage{lmodern} % Ensure lmodern is loaded for pdflatex
\usepackage{tfrupee} % Include tfrupee package

\setlength{\headheight}{1cm} % Set the height of the header box
\setlength{\headsep}{0mm}     % Set the distance between the header box and the top of the text

\usepackage{gvv-book}
\usepackage{gvv}
\usepackage{cite}
\usepackage{physics}
\usepackage{amsmath,amssymb,amsfonts,amsthm}
\usepackage{algorithmic}
\usepackage{graphicx}
\usepackage{textcomp}
\usepackage{xcolor}
\usepackage{txfonts}
\usepackage{listings}
\usepackage{enumitem}
\usepackage{mathtools}
\usepackage{gensymb}
\usepackage{comment}
\usepackage[breaklinks=true]{hyperref}
\usepackage{tkz-euclide} 
\usepackage{listings}
% \usepackage{gvv}                                        
\def\inputGnumericTable{}                                 
\usepackage[latin1]{inputenc}                                
\usepackage{color}                                            
\usepackage{array}                                            
\usepackage{longtable}                                       
\usepackage{calc}                                             
\usepackage{multirow}                                         
\usepackage{hhline}                                           
\usepackage{ifthen}                                           
\usepackage{lscape}
\begin{document}

\bibliographystyle{IEEEtran}

\title{2.8.1}
\author{EE25BTECH11022 - sankeerthan}
% \maketitle
% \newpage
% \bigskip
{\let\newpage\relax\maketitle}

\renewcommand{\thefigure}{\theenumi}
\renewcommand{\thetable}{\theenumi}
\setlength{\intextsep}{10pt} % Space between text and floats


\numberwithin{align}{enumi}
\numberwithin{figure}{enumi}
\renewcommand{\thetable}{\theenumi}

\textbf{Question}: If \brak{-4,3} and \brak{4,3} are two vertices of an equilateral triangle. Find the coordinates of the third vertex, given that the origin lies in the interior of the triangle.

\textbf{solution}: Let 
\begin{align}

\vec{A}= \myvec{-4 \\ 3}, \vec{B}= \myvec{4 \\ 3}
\end{align}

Consider rotating  A  about  B  by 60\degree .
Let the rotation matrix be:
\begin{align}
R_{60\degree} =
\myvec{
\cos 60\degree & -\sin 60\degree \\
\sin 60\degree & \cos 60\degree
}
=
\myvec{\frac{1}{2} & -\frac{\sqrt{3}}{2} \\
\frac{\sqrt{3}}{2} & \frac{1}{2} }
\end{align}

Translate  B  to the origin:
\begin{align}
\vec{A'} = \vec{A} - \vec{B} = \myvec{-4-4 \\ 3-3} = \myvec{-8 \\ 0}
\end{align}

Translate back:
\begin{align}
\vec{C}&= \vec{B}+ R_{60\degree}(\vec{A}) \\ 
       &= \myvec{4 \\  3} +\myvec{\frac{1}{2} & -\frac{\sqrt{3}}{2} \\
\frac{\sqrt{3}}{2} & \frac{1}{2} }\myvec{-8 \\ 0} \\
       &= \myvec{4 \\3}+\myvec{\frac{1}{2}\times(-8) + \left(-\frac{\sqrt{3}}{2}\right)\times0 \\ \frac{\sqrt{3}}{2}\times(-8) +\frac{1}{2}\times 0}\\
       &= \myvec{4 \\3}+\myvec{-4 \\ -4\sqrt{3}} \\
       &=\myvec{4+(-4) \\3+(-4\sqrt{3})} \\
       &= \myvec{0 \\ 3-4\sqrt{3}}
\end{align}

\begin{figure}[!ht]
    \centering
    \includegraphics[width=\linewidth]{figs/points.png}
    \caption{}
\end{figure}
\end{document}